% !TEX TS-program = pdflatex
% !TEX encoding = UTF-8 Unicode

% This file is a template using the "beamer" package to create slides for a talk or presentation
% - Giving a talk on some subject.
% - The talk is between 15min and 45min long.
% - Style is ornate.

% MODIFIED by Jonathan Kew, 2008-07-06
% The header comments and encoding in this file were modified for inclusion with TeXworks.
% The content is otherwise unchanged from the original distributed with the beamer package.

\documentclass{beamer}
\setbeamercovered{invisible}



% Copyright 2004 by Till Tantau <tantau@users.sourceforge.net>.
%
% In principle, this file can be redistributed and/or modified under
% the terms of the GNU Public License, version 2.
%
% However, this file is supposed to be a template to be modified
% for your own needs. For this reason, if you use this file as a
% template and not specifically distribute it as part of a another
% package/program, I grant the extra permission to freely copy and
% modify this file as you see fit and even to delete this copyright
% notice. 


\mode<presentation>
{
  \usetheme{Warsaw}
  % or ...

  % or whatever (possibly just delete it)
}


\usepackage[english]{babel}
% or whatever

\usepackage[utf8]{inputenc}
\usepackage{qrcode}
\usepackage{hyperref}
\usepackage{ytableau}
\usepackage{longtable}
% or whatever

\usepackage{times}
\usepackage[T1]{fontenc}
\usepackage{tikz}
% Or whatever. Note that the encoding and the font should match. If T1
% does not look nice, try deleting the line with the fontenc.

%%% MATH RELATED
\renewcommand{\arraystretch}{1.5}

\title{MATH 180 Strategies of Problem Solving} 
\subtitle
{} % (optional)

\author[W.R. Casper] % (optional, use only with lots of authors)
{W.R. Casper}
% - Use the \inst{?} command only if the authors have different
%   affiliation.

\institute[California State University Fullerton] % (optional, but mostly needed)
{
  Department of Mathematics\\
  California State University Fullerton}
% - Use the \inst command only if there are several affiliations.
% - Keep it simple, no one is interested in your street address.

\subject{Talks}
% This is only inserted into the PDF information catalog. Can be left
% out. 



% If you have a file called "university-logo-filename.xxx", where xxx
% is a graphic format that can be processed by latex or pdflatex,
% resp., then you can add a logo as follows:

% \pgfdeclareimage[height=0.5cm]{university-logo}{university-logo-filename}
% \logo{\pgfuseimage{university-logo}}



% Delete this, if you do not want the table of contents to pop up at
% the beginning of each subsection:
\AtBeginSubsection[]
{
  \begin{frame}<beamer>{Outline}
    \tableofcontents[currentsection,currentsubsection]
  \end{frame}
}


% If you wish to uncover everything in a step-wise fashion, uncomment
% the following command: 

%\beamerdefaultoverlayspecification{<+->}


\begin{document}

\begin{frame}
  \titlepage
\end{frame}

\begin{frame}{Outline}
  \tableofcontents
  % You might wish to add the option [pausesections]
\end{frame}

% Since this a solution template for a generic talk, very little can
% be said about how it should be structured. However, the talk length
% of between 15min and 45min and the theme suggest that you stick to
% the following rules:  

% - Exactly two or three sections (other than the summary).
% - At *most* three subsections per section.
% - Talk about 30s to 2min per frame. So there should be between about
%   15 and 30 frames, all told.

\section{SoPS Lecture 15}

\begin{frame}{Review}
  \begin{block}{Last time}
    \begin{itemize}
\pause
      \item Ready set bet
\pause
      \item Expectation values
    \end{itemize}
  \end{block}
\end{frame}

\subsection{Horsing around}

\begin{frame}{Expected horse position}
  \begin{block}{Problem}
    What's the expected position of each horse after one roll of the dice?
    \begin{itemize}
\pause
      \item sample space: outcomes of rolling two dice
    $$\Omega = \{(1,1),(1,2),(1,3),\dots, (6,6)\}$$
\pause
      \item each horse has a different evaluation function 
\pause
      \item horse $2/3$:
        $$f_{2/3}((j,k)) = \left\lbrace\begin{array}{cc}1, & j+k=2\ \text{or}\ 3\\0, & \text{otherwise}\end{array}\right.$$
    \end{itemize}
  \end{block}
\end{frame}

\begin{frame}{Expected horse position}
  \begin{block}{Horse $2/3$ expected position}
    $$f_{2/3}((j,k)) = \left\lbrace\begin{array}{cc}1, & j+k=2\ \text{or}\ 3\\0, & \text{otherwise}\end{array}\right.$$
    \begin{align*}
      \mathbb{E}(f_{2/3})
        & = \frac{f_{2/3}((1,1))+f_{2/3}((1,2))+\dots+f_{2/3}((6,6))}{|\Omega|}
        & = \frac{3}{36}.
    \end{align*}
  \end{block}
\end{frame}

\begin{frame}{Expected horse position}
  \begin{block}{Horse $2/3$ expected position}
    $$f_{4}((j,k)) = \left\lbrace\begin{array}{cc}1, & j+k=4\\ & \text{otherwise}\end{array}\right.$$
    \begin{align*}
      \mathbb{E}(f_{4})
        & = \frac{f_{4}((1,1))+f_{4}((1,2))+\dots+f_{4}((6,6))}{|\Omega|}
        & = \frac{3}{36}.
    \end{align*}
  \end{block}
\end{frame}

\begin{frame}{Expected horse position}
  \begin{block}{Horse $2/3$ expected position}
    $$f_{5}((j,k)) = \left\lbrace\begin{array}{cc}1, & j+k=5\\ & \text{otherwise}\end{array}\right.$$
    \begin{align*}
      \mathbb{E}(f_{5})
        & = \frac{f_{5}((1,1))+f_{5}((1,2))+\dots+f_{5}((6,6))}{|\Omega|}
        & = \frac{4}{36}.
    \end{align*}
  \end{block}
\end{frame}

\begin{frame}{Expected horse position}
  \begin{block}{Horse $2/3$ expected position}
    $$f_{6}((j,k)) = \left\lbrace\begin{array}{cc}1, & j+k=6\\ & \text{otherwise}\end{array}\right.$$
    \begin{align*}
      \mathbb{E}(f_{6})
        & = \frac{f_{6}((1,1))+f_{6}((1,2))+\dots+f_{6}((6,6))}{|\Omega|}
        & = \frac{5}{36}.
    \end{align*}
  \end{block}
\end{frame}

\begin{frame}{Expected horse position}
  \begin{block}{Horse $2/3$ expected position}
    $$f_{7}((j,k)) = \left\lbrace\begin{array}{cc}1, & j+k=7\\ & \text{otherwise}\end{array}\right.$$
    \begin{align*}
      \mathbb{E}(f_{7})
        & = \frac{f_{7}((1,1))+f_{7}((1,2))+\dots+f_{7}((6,6))}{|\Omega|}
        & = \frac{6}{36}.
    \end{align*}
  \end{block}
\end{frame}

\begin{frame}{Expected horse position}
  \begin{block}{Problem}
    What's the expected position of each horse after one roll of the dice?
  \end{block}
  \textbf{Solution:}
  \begin{center}
    \begin{tabular}{l|c|c|c|c|c|c|c|c|c|}
      horse    & 2/3 & 4 & 5 & 6 & 7 & 8 & 9 & 10 & 11/12\\\hline
      exp. pos & $\frac{1}{12}$ & $\frac{1}{12}$ & $\frac{1}{9}$ & $\frac{5}{36}$ & $\frac{1}{6}$ & $\frac{5}{36}$ & $\frac{1}{9}$ & $\frac{1}{12}$ & $\frac{1}{12}$
    \end{tabular}
  \end{center}
\end{frame}

\begin{frame}{Expected horse position}
  \begin{block}{Problem}
    What's the expected position of each horse after two rolls of the dice?
  \end{block}
    \begin{itemize}
\pause
      \item sample space: outcomes of rolling two dice
    $$\Omega = \{(1,1,1,1),(1,1,1,2),(1,1,1,3),\dots, (6,6,6,6)\}$$
\pause
      \item size of sample space:
    $$|\Omega| = 6^4 = 1296$$
    \end{itemize}
\pause
    $${\color{red}\text{sample space getting bigger}\dots}$$
\end{frame}

\begin{frame}{Non-uniform probabilty}
    \begin{itemize}
\pause
     \item alternative sample space:
  $$\Omega = \{(2/3,2/3),(2/3,4),(2/3, 5)\dots, (11/12,10), (11/12,11/12)\}$$
\pause
      \item size of sample space:
    $$|\Omega| = 9^2 = 81$$
\pause
      \item tradeoff is probabilities are no longer uniform!
  \end{itemize}
\end{frame}

\begin{frame}{Non-uniform probabilty}
  \begin{block}{Problem}
    What is the probabilty of $(2/3,2/3)$?
  \end{block}
\pause
  \begin{block}{Theorem}
    If $A$ and $B$ are independent events, then
    $$\mathbb{P}(\text{$A$ and $B$}) = \mathbb{P}(A)\mathbb{P}(B).$$
  \end{block}
\pause
  $$\mathbb P((2/3,2/3)) = \mathbb P(2/3)\mathbb P(2/3) = \frac{3}{36}\cdot\frac{3}{36} = \frac{1}{144}.$$
\end{frame}

\begin{frame}{Non-uniform probabilty}
  \begin{block}{Problem}
    What is the probabilty of $(4,7)$?
  \end{block}
\pause
  $$\mathbb P((4,7)) = \mathbb P(4)\mathbb P(7) = \frac{3}{36}\cdot\frac{6}{36} = \frac{1}{72}.$$
\end{frame}

\begin{frame}{Probability table}
  $$\text{probabilities (divided by $1296$)}$$
{\tiny
\begin{center}
  \begin{tabular}{l|c|c|c|c|c|c|c|c|c}
  $\mathbb{P}$ & 2/3 & 4 & 5 & 6 & 7 & 8 & 9 & 10 & 11/12\\\hline
           2/3 & 9 & 9 & 12 & 15 & 18 & 15 & 12 & 9 & 9\\\hline
            4  & 9 & 9 & 12 & 15 & 18 & 15 & 12 & 9 & 9\\\hline
            5  & 12 & 12 & 16 & 20 & 24 & 20 & 16 & 12 & 12\\ \hline
            6  & 15 & 15 & 20 & 25 & 30 & 25 & 20 & 15 & 15\\\hline
            7  & 18 & 18 & 24 & 30 & 36 & 30 & 24 & 18 & 18\\\hline
            8  & 15 & 15 & 20 & 25 & 30 & 25 & 20 & 15 & 15\\\hline
            9  & 12 & 12 & 16 & 20 & 24 & 20 & 16 & 12 & 12\\\hline
           10  & 9 & 9 & 12 & 15 & 18 & 15 & 12 & 9 & 9\\\hline
          11/12& 9 & 9 & 12 & 15 & 18 & 15 & 12 & 9 & 9
  \end{tabular}
\end{center}
}
\end{frame}

\begin{frame}{Non-uniform probabilities}
\begin{itemize}
\pause
\item sample space
$$\Omega = \{o_1,o_2,\dots, o_N\}$$
\pause
\item probability of $o_i$ is $p_i=\mathbb P(o_i)$
\pause
\item probability of event $E = \{o_{j_1},o_{j_2},\dots, o_{j_r}\}$ is
$$\mathbb P(E) = p_{j_1} + p_{j_2} + \dots + p_{j_r}.$$
\pause
\item expectation of $f: \Omega\rightarrow\mathbb R$
$$\mathbb E(f) = p_1 f(o_1) + p_2 f(o_2) + p_3 f(o_3) + \dots + p_N f(o_N).$$
\end{itemize}
\end{frame}

\begin{frame}{Expected horse position}
  \begin{block}{Problem}
    What is the expected position of the $2/3$ horse after two rolls?
  \end{block}
\pause
{\tiny
  \begin{align*}
    \mathbb{E}(\text{2/3 pos})
    & = 5\cdot \mathbb P((2/3,2/3))\quad \text{2 moves + 3 bonus}\\
    & + 1\cdot \mathbb P((2/3,4)) + 1\cdot\mathbb P((2/3,5)) + \dots + 1\cdot\mathbb P((2/3,11/12))\\
    & + 1\cdot \mathbb P((4,2/3)) + 1\cdot\mathbb P((5,2/3)) + \dots + 1\cdot\mathbb P((11/12,2/3))\\
    & = 5\cdot \frac{9}{1296}\\
    & + \frac{9}{1296} + \frac{12}{1296} + \frac{15}{1296} + \frac{18}{1296} +  \frac{15}{1296} +  \frac{12}{1296} + \frac{9}{1296} + \frac{9}{1296} \\
    & + \frac{9}{1296} + \frac{12}{1296} + \frac{15}{1296} + \frac{18}{1296} +  \frac{15}{1296} +  \frac{12}{1296} + \frac{9}{1296} + \frac{9}{1296} \\
    & = \frac{45 + 54 + 48 + 60 + 36}{1296} = \frac{243}{1296} = 0.1875
  \end{align*}
}
\end{frame}

\begin{frame}{Expected horse position}
  \begin{block}{Problem}
    What is the expected position of each horse after two rolls?
  \end{block}
\pause
  \textbf{Solution:}
  {\tiny
  \begin{center}
    \begin{tabular}{l|c|c|c|c|c|c|c|c|c|}
      horse    & 2/3 & 4 & 5 & 6 & 7 & 8 & 9 & 10 & 11/12\\\hline
      exp. pos & $0.1875$ & $0.1875$ & $0.2469$ & $0.2971$ & $0.3333$ & $0.2971$ & $0.2469$ & $0.1875$ & $0.1875$
    \end{tabular}
  \end{center}
  }
\end{frame}
\subsection{Wrapping up}

\begin{frame}{Final thoughts}
  \begin{block}{Reflection}
    \begin{itemize}
\pause
      \item Suppose that there is no bonus for horses anymore.  Work out by hand how this changes the expected value of the 2/3 horse after two dice rolls.
    \end{itemize}
  \end{block}
\begin{center}
Submit your responses by \textbf{Tonight!}
\end{center}
\end{frame}


\end{document}


